% ============================================================================
% PART II: PHYSICS ORGANIZED BY OBSERVABLE FAMILY
% Section 7: Defects, Caveats, and Outlook
% ============================================================================
%
% Purpose
%   Absorbs the rest of the observable taxonomy:
%     - Monopoles and line operators (defect / global-structure family).
%     - The unified "what is strictly topological vs IR-effective vs
%       theory-side vs model-dependent" discussion (subsuming the localized
%       caveat subsections of Secs. 5.6 and 6.4).
%     - Guide du routard: which subfields use which observables, which papers
%       to read next, where to go for formal vs phenomenological questions.
%
% Status: NEW WRITING, plus EXTRACT of any relevant material from
%   archive_2026-05-01/anomalies_boundaries_topological_response.tex (for
%   'iffiness' synthesis) and Sec. 3 (for monopole motivation).
%
% Length target: ~12--15 pages total
%
% Subsection plan (from Final QFT Project Overview.md)
%   7.1  Monopoles, line operators, and global structure
%   7.2  When the observables are really topological, and when they are only effective
%   7.3  Outlook and guide du routard
%
% References: monopoles and defects
%   Dirac 1931 (\cite{Dirac1931QuantisedSingularities})
%   't Hooft 1974 (\cite{tHooft1974MagneticMonopoles})
%   Polyakov 1974 (\cite{Polyakov1974ParticleSpectrum})
%   Manton-Sutcliffe (\cite{MantonSutcliffeTopologicalSolitons})
%   Shnir (\cite{ShnirMagneticMonopoles})
%   Weinberg (\cite{WeinbergClassicalSolutionsQFT})
%   Rajantie 2024 (\cite{Rajantie2024MonopoleTheoryOverview})
%   Fairbairn et al. 2021 (\cite{FairbairnMonopolesRevisited})
%   PDG Monopoles (\cite{PDGMagneticMonopolesPlaceholder})
%
% References: line operators / global form / generalized symmetries
%   Witten 1989 (\cite{Witten1989Jones})
%   Gaiotto-Kapustin-Seiberg-Willett 2015 (\cite{GaiottoKapustinSeibergWillett2015GeneralizedSymmetries})
%   Kapustin-Saulina line operators (\cite{KapustinSaulinaLineOperatorsPlaceholder})
%   Kapustin-Witten 2007 (\cite{KapustinWitten2007ElectricMagneticLanglands})
%     -- advanced outlook pointer only
%
% References: caveats / organizing principle
%   Witten 2016 (\cite{Witten2016ThreeLecturesTopologicalPhases})
%   Stern 2008 (\cite{Stern2008AnyonsQHEReview})
%   Birmingham et al. 1991 (\cite{BirminghamBlauRakowskiThompson1991TopologicalFieldTheory})
%   Wen 1995 (\cite{Wen1995TopologicalOrdersEdgeExcitations})
%
% Writing notes
%   - 7.1 covers the defect family in one place: why monopole sectors are
%     topological, why the underlying theory is not, and how line operators
%     probe the global form of the gauge group.
%   - 7.2 is the final, unified version of the "iffiness" theme: four
%     categories of observables --- strict TQFT, IR response, theory-side
%     topological, model-dependent --- per
%     tqft_observables_literature_review.md Sec. 9.
%   - 7.3 is a short guide-du-routard pointing the reader to the appropriate
%     next references depending on their interest.

\section{Defects, Caveats, and Outlook}
\label{sec:defects-caveats-outlook}

% TODO: Short preamble stating that the remaining observable families
% (defects, global structure) and the "how real is this observable" question
% live here.

\subsection{Monopoles, Line Operators, and Global Structure}
\label{subsec:monopoles-line-operators}

% TODO: Dirac monopoles and charge quantization as a topological consistency
% condition; 't Hooft-Polyakov monopoles as finite-energy solitons with
% maps S^2_infty -> G/H (use pi_2(G/H) from Sec. 2.5). Contrast singular
% Dirac vs smooth 't Hooft-Polyakov. Magnetic charge sector != topological
% theory. Finally: Wilson / 't Hooft lines as probes of the global form of
% the gauge group (e.g. SU(2) vs SO(3)) and preview of generalized
% global symmetries.
% Cite: \cite{Dirac1931QuantisedSingularities}, \cite{tHooft1974MagneticMonopoles},
%       \cite{Polyakov1974ParticleSpectrum},
%       \cite{MantonSutcliffeTopologicalSolitons},
%       \cite{ShnirMagneticMonopoles}, \cite{WeinbergClassicalSolutionsQFT},
%       \cite{Rajantie2024MonopoleTheoryOverview},
%       \cite{FairbairnMonopolesRevisited},
%       \cite{PDGMagneticMonopolesPlaceholder},
%       \cite{GaiottoKapustinSeibergWillett2015GeneralizedSymmetries},
%       \cite{KapustinSaulinaLineOperatorsPlaceholder}

\subsection{When the Observables Are Really Topological, and When They Are Only Effective}
\label{subsec:when-topological}

% From tqft_observables_literature_review.md Sec. 9: four-way organizing
% table for the observables discussed throughout Part II.
%
%   (a) Strict TQFT observables
%         Wilson-line linking, knot data, Hilbert spaces on closed surfaces.
%         Metric-independent in the exact theory.
%         Examples in this paper: Sec. 4.4--4.5.
%
%   (b) IR topological response observables
%         Hall conductance, edge anomaly coefficient, GSD on topology.
%         Universal only below the gap and after decoupling assumptions.
%         Examples: Sec. 5.1--5.3.
%
%   (c) Topology-controlled but not directly measurable theory observables
%         Instanton number, chi_t.
%         Defined cleanly in Euclidean QFT / lattice QCD; enter observables
%         indirectly (eta', axion mass, theta-dependence of E).
%         Examples: Sec. 6.1--6.2.
%
%   (d) Topology-inspired but model-dependent observables
%         Axion couplings in specific UV completions, monopole production
%         rates, interferometric signatures in mesoscopic devices.
%         Examples: Sec. 5.6, 6.3, 7.1.
%
% This subsection is where the lecture-flavor "iffiness" lives in its most
% structural form, after the local caveat subsections in Secs. 5 and 6.
% Cite: \cite{Witten2016ThreeLecturesTopologicalPhases}, \cite{Stern2008AnyonsQHEReview},
%       \cite{Wen1995TopologicalOrdersEdgeExcitations},
%       \cite{BirminghamBlauRakowskiThompson1991TopologicalFieldTheory}

\subsection{Guide du Routard: Where to Go Next}
\label{subsec:guide-du-routard}

% Short pointer list (tqft_observables_literature_review.md Sec. 8):
%   - Response and experiment: \cite{Wen1995TopologicalOrdersEdgeExcitations},
%     \cite{Stern2008AnyonsQHEReview},
%     \cite{HalperinJain2020FQHENewDevelopments},
%     anyon-interferometry experiments.
%   - Formal TQFT: \cite{Atiyah1988TQFT}, \cite{Witten1988TQFT},
%     \cite{Witten1989Jones},
%     \cite{BirminghamBlauRakowskiThompson1991TopologicalFieldTheory}.
%   - QCD topology / axions: \cite{Dine2000TASIStrongCP},
%     \cite{Hook2018TASIStrongCPAxions}, \cite{Marsh2016AxionCosmology},
%     \cite{Witten1979U1GoldstoneBoson}, \cite{Veneziano1979U1WithoutInstantons},
%     lattice literature.
%   - Monopoles / solitons / defects: \cite{Dirac1931QuantisedSingularities},
%     \cite{tHooft1974MagneticMonopoles}, \cite{Polyakov1974ParticleSpectrum},
%     \cite{MantonSutcliffeTopologicalSolitons}, \cite{ShnirMagneticMonopoles},
%     \cite{Rajantie2024MonopoleTheoryOverview},
%     \cite{FairbairnMonopolesRevisited}.
%   - Modern operator language and generalized symmetries:
%     \cite{GaiottoKapustinSeibergWillett2015GeneralizedSymmetries},
%     later work (Kapustin-Saulina placeholder).
%
% Close with a final paragraph: TQFT reorganizes gauge theory around global
% data; makes protected observables transparent (Hall conductance, chi_t,
% Dirac quantization); exposes anomaly/boundary structure structurally
% (inflow -> edge modes); and provides a computable language for
% universality classes across condensed matter and particle physics.
